\chapter*{Abstract}
\addcontentsline{toc}{chapter}{Abstract}

European companies run much of their digital work on American cloud services, and
American law allows the American government to demand that data even when it is
stored in Europe. This thesis asks whether European cloud providers can offer a
real alternative, using Alps Alpine as the example. Eighteen European providers were
examined in four studies: their sovereignty, judged with the European Commission's
Cloud Sovereignty Framework; the services they offer, compared with the sixty
\acs{AWS} services Alps Alpine uses or expects to need; the money behind them; and
the speed and price of their machines, measured in standard tests.

Five of the eighteen providers meet every requirement set in this study, but none
reaches full sovereignty, because no European provider runs on chips or firmware made
in Europe. No European provider can yet replace \acs{AWS} on its own: the broadest
covers a little under three quarters of the services Alps Alpine uses, and most cover
around a quarter, with artificial intelligence the weakest area. On machines of the
same size, however, European providers are as fast as \acs{AWS} and often much
cheaper. The providers that offer the most control over data tend to offer the
fewest services, although STACKIT shows that a provider can grow without giving up
control.

The thesis concludes with a recommendation matrix that matches providers to the
sensitivity of the data and the kind of service needed. For Alps Alpine it
recommends moving in steps: confidential data first, to STACKIT, then internal work
to T-Cloud Public or OVHcloud, and artificial intelligence last, once running
openly available models at a sovereign provider has been tested.

\chapter*{Zusammenfassung}
\addcontentsline{toc}{chapter}{Zusammenfassung}

Europäische Unternehmen betreiben einen großen Teil ihrer digitalen Arbeit bei
amerikanischen Cloud-Diensten, und das amerikanische Recht erlaubt der
amerikanischen Regierung, diese Daten anzufordern, selbst wenn sie in Europa
gespeichert sind. Diese Arbeit untersucht am Beispiel von Alps Alpine, ob
europäische Cloud-Anbieter eine echte Alternative bieten können. Achtzehn europäische
Anbieter wurden in vier Studien betrachtet: ihre Souveränität, bewertet mit dem Cloud
Sovereignty Framework der Europäischen Kommission; die angebotenen Dienste, verglichen
mit den sechzig \acs{AWS}-Diensten, die Alps Alpine nutzt oder künftig benötigt; das
Geld hinter den Anbietern; sowie Geschwindigkeit und Preis ihrer Maschinen, gemessen
mit Standardtests.

Fünf der achtzehn Anbieter erfüllen alle in dieser Arbeit festgelegten
Anforderungen, doch keiner erreicht volle Souveränität, da kein europäischer Anbieter
mit in Europa hergestellten Chips oder Firmware arbeitet. Kein europäischer Anbieter
kann \acs{AWS} bisher allein ersetzen: Der breiteste deckt knapp drei Viertel der von
Alps Alpine genutzten Dienste ab, die meisten etwa ein Viertel, wobei künstliche
Intelligenz der schwächste Bereich ist. Auf gleich großen Maschinen sind europäische
Anbieter jedoch ebenso schnell wie \acs{AWS} und oft deutlich günstiger. Die Anbieter
mit der größten Kontrolle über die Daten bieten tendenziell die wenigsten Dienste an,
auch wenn STACKIT zeigt, dass ein Anbieter wachsen kann, ohne die Kontrolle
aufzugeben.

Die Arbeit schließt mit einer Empfehlungsmatrix, die Anbieter der Sensibilität der
Daten und der Art des benötigten Dienstes zuordnet. Für Alps Alpine empfiehlt sie ein
schrittweises Vorgehen: zuerst vertrauliche Daten zu STACKIT, dann interne Anwendungen
zu T-Cloud Public oder OVHcloud und künstliche Intelligenz zuletzt, sobald der Betrieb
frei verfügbarer Modelle bei einem souveränen Anbieter erprobt ist.
